% =====================================================================
\BookChapter[The Wolfram Language]{The Wolfram Language: Computing with Symbols}
\label{ch:wolfram}

\begin{researchquestion}
When should the computer manipulate the formula itself and not just numbers?
\end{researchquestion}
\OrientationNote
% =====================================================================

\begin{factsheet}{Wolfram Language}
First appeared & 1988 as the language of Mathematica 1.0; named ``Wolfram Language'' in 2013; Mathematica 14 (2024)\\
Designed by & Stephen Wolfram and Wolfram Research\\
Paradigms & symbolic, functional, rule-based (term rewriting), with procedural constructs\\
Typing & dynamic; everything is a symbolic \emph{expression}\\
Execution & interpreted by the Wolfram kernel; numerical routines compiled\\
File extensions & \file{.nb} (notebook), \file{.wl} and \file{.m} (packages and scripts)\\
Cost & commercial; many universities have site licences. The Wolfram Engine is free for
  developers, and Wolfram|Alpha uses the language on the web\\
Used for & symbolic mathematics (algebra, calculus, differential equations), deriving and
  checking formulas, teaching, physics, control theory, rapid prototyping, interactive
  demonstrations\\
Famous software & Mathematica, Wolfram|Alpha, the Wolfram Demonstrations Project, System
  Modeler (Modelica)\\
\end{factsheet}

All the languages so far compute with \emph{numbers}. Ask \py for the derivative of
$x^3 e^{-x}$ and it cannot help, unless you install a special library. The Wolfram
Language computes with \emph{symbols}: it can expand, factor, differentiate and integrate
formulas, solve equations exactly, and simplify expressions --- the things you do with pen
and paper in a mathematics or mechanics course. For an engineer, a computer algebra system
is the tool for deriving a formula, checking a hand derivation, or finding the exact
solution of a model before writing a numerical code for the general case.

\LearningOutcomes{Read Wolfram Language syntax; expand, solve, differentiate and integrate expressions symbolically; solve a differential equation exactly and numerically; recognise functional and rule-based programming; attach units to quantities; and check a symbolic result before relying on it.}

\section{History and philosophy}

\begin{history}
Stephen Wolfram, a British physicist who received his PhD from Caltech at the age of 20,
wrote a computer algebra system called SMP in 1979--81 for his particle-physics work.
Dissatisfied with it, he founded Wolfram Research and released Mathematica in 1988. Its
notebook interface --- documents that mix text, formulas, code, results and graphics ---
was revolutionary and directly inspired today's Jupyter notebooks. In 2009 the company
launched Wolfram|Alpha, a ``computational knowledge engine'' that answers questions typed
in plain English using the same language underneath, and in 2013 the language itself was
given its own name.
\end{history}

The language rests on one idea: \emph{everything is an expression}. A number, a formula, a
list, a graphic, a piece of code and a document are all expressions of the form
\code{head[arg1, arg2, \ldots]}. For example, $x + y^2$ is internally
\code{Plus[x, Power[y, 2]]}. Computation means transforming expressions by applying
rules, until no more rules apply. This makes symbolic and numeric work, programs and data
all uniform. A second principle is that the language should contain as much built-in
knowledge as possible: it has more than 6000 built-in functions and curated data about
chemical elements, materials, units, physical constants, geography and much more.

\section{Getting started}

Mathematica runs on Windows, macOS and Linux, and in a web browser (Wolfram Cloud).
Students at many universities get a licence through their institution. The free
\textbf{Wolfram Engine} runs the language from scripts, Jupyter notebooks or \py, and
\textbf{Wolfram|Alpha} (\url{https://www.wolframalpha.com}) accepts both English and
Wolfram Language input.

In a notebook you type an input and press \code{Shift+Enter}. Inputs are labelled
\code{In[n]:=} and results \code{Out[n]=}:
\begin{lstlisting}[language=Mathematica]
In[1]:= Print["Hello, engineers!"]
        Hello, engineers!
In[2]:= 2 + 3
Out[2]= 5
In[3]:= 1/3 + 1/4
Out[3]= 7/12
In[4]:= N[%]          (* % is the last result; N gives a numerical value *)
Out[4]= 0.583333
In[5]:= N[Pi, 30]
Out[5]= 3.14159265358979323846264338328
\end{lstlisting}
Notice that \code{1/3 + 1/4} gives the \emph{exact} fraction $7/12$: integers and fractions
are kept exact unless you ask for a numerical approximation with \code{N} or include a
decimal point (\code{1./3}). \code{N[Pi, 30]} shows that numbers can have arbitrary
precision.

\section{Syntax essentials}

The syntax differs from every other language in these notes, so it is worth listing the
conventions.
\begin{itemize}
\item Built-in functions start with a capital letter and are written in full:
  \code{Sin}, \code{Exp}, \code{Integrate}, \code{Solve}, \code{Plot}. To avoid clashes,
  give your own names a lower-case first letter.
\item Function arguments go in \textbf{square brackets}: \code{Sin[x]}, not \code{Sin(x)}.
  Round brackets are only for grouping.
\item \textbf{Curly braces} make lists: \code{\{1, 2, 3\}}. Vectors are lists, and a
  matrix is a list of lists. Parts are extracted with double square brackets:
  \code{list[[2]]}. Indices start at 1.
\item Multiplication can be written with a space: \code{2 x y} means $2xy$.
  \code{x\^{}2} is a power.
\item \code{=} assigns, \code{:=} defines a rule that is evaluated each time it is used,
  \code{==} is an equation, and \code{->} is a replacement rule.
\item A semicolon at the end of an input suppresses the output. Comments are written
  \code{(* like this *)}.
\item Built-in constants: \code{Pi}, \code{E}, \code{I} ($\sqrt{-1}$), \code{Infinity},
  \code{Degree}.
\end{itemize}

\section{Symbolic mathematics}

\subsection*{Algebra is done exactly}

\begin{lstlisting}[language=Mathematica]
In[1]:= Expand[(a + b)^3]
Out[1]= a^3 + 3 a^2 b + 3 a b^2 + b^3
In[2]:= Factor[x^2 - 3 x + 2]
Out[2]= (-2 + x) (-1 + x)
In[3]:= Solve[x^2 - 3 x + 2 == 0, x]
Out[3]= {{x -> 1}, {x -> 2}}
In[4]:= Solve[a x^2 + b x + c == 0, x]
Out[4]= {{x -> (-b - Sqrt[b^2 - 4 a c])/(2 a)}, {x -> (-b + Sqrt[b^2 - 4 a c])/(2 a)}}
In[5]:= Simplify[Sin[t]^2 + Cos[t]^2]
Out[5]= 1
\end{lstlisting}
\code{Solve} returns a list of \term{replacement rules}. To use a solution, apply the rule
with \code{/.} (``replace all''):
\begin{lstlisting}[language=Mathematica]
In[6]:= x^3 /. {x -> 2}
Out[6]= 8
In[7]:= vt = v /. Solve[m g == c v^2 && v > 0, v, Reals][[1]]
\end{lstlisting}
The last line asks for the positive speed at which drag balances weight, and the result is
the familiar $\sqrt{gm}/\sqrt{c}$ (with conditions on the signs of the parameters, which the
system reports carefully).

\subsection*{Derivatives, integrals, limits and series are single commands}

\begin{lstlisting}[language=Mathematica]
In[1]:= D[x^3 Exp[-x], x]
Out[1]= 3 E^-x x^2 - E^-x x^3
In[2]:= Integrate[Sin[x]^2, x]
Out[2]= x/2 - 1/4 Sin[2 x]
In[3]:= Integrate[Exp[-x^2], {x, -Infinity, Infinity}]
Out[3]= Sqrt[Pi]
In[4]:= Limit[Sin[x]/x, x -> 0]
Out[4]= 1
In[5]:= Series[Tanh[x], {x, 0, 7}]
Out[5]= x - x^3/3 + (2 x^5)/15 - (17 x^7)/315 + O[x]^8
In[6]:= Sum[1/k^2, {k, 1, Infinity}]
Out[6]= Pi^2/6
\end{lstlisting}

\subsection*{Some differential equations can be solved exactly}

\code{DSolve} finds exact solutions of ordinary differential equations where they exist.
For a damped harmonic oscillator, $\ddot x + 2\zeta\omega\dot x + \omega^2 x = 0$:
\begin{lstlisting}[language=Mathematica]
In[1]:= DSolve[{x''[t] + 2 z w x'[t] + w^2 x[t] == 0, x[0] == 1, x'[0] == 0}, x[t], t]
\end{lstlisting}
returns the general formula in terms of exponentials. For our drag problem,
\begin{lstlisting}[language=Mathematica]
In[2]:= sol = DSolve[{m v'[t] == m g - c v[t]^2, v[0] == 0}, v[t], t];
In[3]:= FullSimplify[v[t] /. sol[[1]], {m > 0, g > 0, c > 0, t > 0}]
\end{lstlisting}
gives an expression equivalent to $\sqrt{mg/c}\,\tanh\!\big(t\sqrt{gc/m}\big)$, which is
equation~\eqref{eq:dragexact} of the opening guide. This is exactly how one would \emph{derive}
the exact solution that all our numerical programs were checked against. When no exact
solution exists, \code{NDSolve} solves the equation numerically.

\subsection*{Matrices of integers give exact answers}

\begin{lstlisting}[language=Mathematica]
In[1]:= A = {{4, -2, 1}, {-2, 4, -2}, {1, -2, 4}};
In[2]:= LinearSolve[A, {11, -16, 17}]
Out[2]= {1, -2, 3}
In[3]:= Eigenvalues[A]
Out[3]= {1/2 (9 + Sqrt[33]), 3, 1/2 (9 - Sqrt[33])}
In[4]:= Det[A]
Out[4]= 36
In[5]:= MatrixForm[Inverse[A]]
\end{lstlisting}
Because the matrix contains integers, the eigenvalues come out \emph{exactly}; MATLAB gave
the same values as decimals (Chapter~\ref{ch:matlab}). Symbolic matrices work too: you can
compute the stiffness matrix of a beam element with symbolic length $L$ and modulus $E$.

\section{Numerics and graphics}

The Wolfram Language is also a capable numerical system:
\begin{lstlisting}[language=Mathematica]
In[1]:= FindRoot[x^3 - 2 x - 5, {x, 2}]
Out[1]= {x -> 2.09455}
In[2]:= NIntegrate[Sin[x], {x, 0, Pi}]
Out[2]= 2.
In[3]:= sol = NDSolve[{v'[t] == 9.81 - 0.25/80 v[t]^2, v[0] == 0}, v, {t, 0, 20}];
In[4]:= v[10] /. sol
Out[4]= {52.7496}
In[5]:= Plot[Evaluate[v[t] /. sol], {t, 0, 20}, Frame -> True,
             FrameLabel -> {"t (s)", "v (m/s)"}]
\end{lstlisting}
\code{Plot}, \code{ListPlot}, \code{Plot3D}, \code{ContourPlot} and \code{ParametricPlot}
produce graphics that appear directly in the notebook. \code{Manipulate} creates an
interactive demonstration with sliders in one line:
\begin{lstlisting}[language=Mathematica]
Manipulate[
 Plot[Sqrt[m 9.81/c] Tanh[t Sqrt[9.81 c/m]], {t, 0, 20}, PlotRange -> {0, 130}],
 {{m, 80}, 40, 120}, {{c, 0.25}, 0.05, 1}]
\end{lstlisting}

\section{Functional and rule-based programming}

Wolfram Language code is typically \emph{functional}: instead of loops that modify
variables, you apply functions to lists.
\begin{lstlisting}[language=Mathematica]
In[1]:= Table[k^2, {k, 1, 5}]
Out[1]= {1, 4, 9, 16, 25}
In[2]:= Map[Sqrt, {1, 4, 9}]          (* shorthand: Sqrt /@ {1, 4, 9} *)
Out[2]= {1, 2, 3}
In[3]:= Select[Range[20], PrimeQ]
Out[3]= {2, 3, 5, 7, 11, 13, 17, 19}
In[4]:= #^2 + 1 & /@ {1, 2, 3}        (* a pure (anonymous) function *)
Out[4]= {2, 5, 10}
In[5]:= Total[{3, 1, 4}]
Out[5]= 8
\end{lstlisting}
In a pure function, \code{\#} stands for the argument and \code{\&} marks the end of the
function, so \code{\#\^{}2 + 1 \&} is the counterpart of \py's \code{lambda x: x**2 + 1}.

User functions are defined by \term{patterns}. The underscore in \code{x\_} means ``match
any expression and call it \code{x}'':
\begin{lstlisting}[language=Mathematica]
terminalVelocity[m_, c_, g_ : 9.81] := Sqrt[m g/c]   (* g_ : 9.81 = default value *)
terminalVelocity[80, 0.25]                           (* 56.0286 *)
terminalVelocity[m, c]                               (* symbolic: 3.13209 Sqrt[m/c] *)
\end{lstlisting}
The same function works with numbers and with symbols. Patterns can also restrict the
type of argument (\code{n\_Integer}) or add conditions, which allows definitions that look
like mathematical case distinctions.

\section{Case study: the falling skydiver}

Here is the Euler method written in the functional style. \code{NestList[f, x0, n]} applies
\code{f} repeatedly and returns the list $\{x_0, f(x_0), f(f(x_0)), \ldots\}$ --- exactly the
sequence of Euler updates:
\begin{lstlisting}[language=Mathematica,style=numbered]
(* Falling object with quadratic drag: explicit Euler method *)
m = 80.; c = 0.25; g = 9.81; dt = 0.5; tEnd = 20.;
nSteps = Round[tEnd/dt];
vt = Sqrt[m g/c];

eulerStep[v_] := v + dt (g - c/m v^2);
vEuler = NestList[eulerStep, 0., nSteps];   (* {v0, v1, ..., v40} *)
times  = Range[0., tEnd, dt];
vExact = vt Tanh[g times/vt];               (* listable: works on the whole list *)

table = Transpose[{times, vEuler, vExact, vEuler - vExact}][[1 ;; -1 ;; 4]];
Print[TableForm[table, TableHeadings -> {None, {"t (s)", "v_Euler", "v_exact", "error"}}]]
Print["Terminal velocity: ", NumberForm[vt, {6, 4}], " m/s"]
\end{lstlisting}
Run in a Mathematica notebook, the program prints:
\outfile[title={\footnotesize\sffamily\color{muted}Output (Mathematica notebook)}]{wolfram/drag.out}
The numbers agree with the tables of the other chapters to all the digits shown. Only
the appearance differs: by default the Wolfram Language shows six significant digits,
drops trailing zeros, and writes a whole number such as 2 as \code{2.} to mark it as a
decimal (not exact) value. The table is wrapped in
\code{Print} so that it also appears when the file is run as a script with
\code{wolframscript}; in a notebook, \code{TableForm} alone would be displayed. Note \code{[[1 ;; -1 ;; 4]]}, the
Wolfram Language's slice notation (start, end, step), and that the arithmetic on
\code{times} is applied element-wise automatically because \code{Tanh} and \code{*} are
\emph{listable}. A procedural version with \code{Do} or \code{For} loops is also possible,
but is considered poor style.

\section{Units and data}

The language has built-in physical units and curated data, which helps to avoid the
Mars-Climate-Orbiter type of error (Chapter~\ref{ch:howcoderuns}):
\begin{lstlisting}[language=Mathematica]
In[1]:= N[UnitConvert[Quantity[60, "Miles"/"Hours"], "Meters"/"Seconds"]]
Out[1]= 26.8224 m/s
In[2]:= Quantity[80, "Kilograms"] Quantity[9.81, "Meters"/"Seconds"^2]
Out[2]= 784.8 kg m/s^2
In[3]:= UnitConvert[%, "Newtons"]
Out[3]= 784.8 N
ElementData["Iron", "Density"]              (* density of iron, with units *)
Entity["Element", "Copper"]["ThermalConductivity"]
\end{lstlisting}

\section{Computer algebra in other languages}

You do not need Mathematica to do symbolic mathematics:
\begin{itemize}
\item \textbf{SymPy} is a free \py library for symbolic mathematics. It covers much of the
  same ground (\code{diff}, \code{integrate}, \code{solve}, \code{dsolve}, \code{series})
  and can generate C, Fortran or NumPy code from formulas:
\begin{lstlisting}[language=Python]
import sympy as sp
x = sp.symbols("x")
sp.diff(x**3 * sp.exp(-x), x)          # -x**3*exp(-x) + 3*x**2*exp(-x)
sp.integrate(sp.sin(x)**2, x)          # x/2 - sin(x)*cos(x)/2
\end{lstlisting}
\item \textbf{Maple} is a commercial computer algebra system popular in engineering
  education and in some industries; \textbf{Maxima} is a free system with a long history.
\item MATLAB's \textbf{Symbolic Math Toolbox} provides symbolic computation inside
  MATLAB.
\item Many graphing calculators include a computer algebra system.
\end{itemize}

\begin{engnote}[Engineering note: trust, but verify]
Computer algebra is powerful but not infallible. Results may be correct only under
assumptions you did not state (is $\sqrt{x^2} = x$? Only if $x \ge 0$), may be returned in
an unrecognisable form, or may take forever. Always check a symbolic result: substitute
numbers and compare with a numerical calculation, check limiting cases, and check the
units.
\end{engnote}

\section{Strengths and weaknesses}

\begin{center}
\small
\begin{tabularx}{\linewidth}{@{}>{\raggedright\arraybackslash}X>{\raggedright\arraybackslash}X@{}}
\toprule
\textbf{Strengths} & \textbf{Weaknesses}\\
\midrule
World-class symbolic mathematics & Commercial and proprietary\\
Exact and arbitrary-precision arithmetic & Unusual syntax; code is hard for non-users to read\\
Excellent notebooks, graphics and interactivity & Slower than compiled languages for large numerical loops\\
Huge built-in function library with units and curated data & Small community outside science and education\\
Very concise code & Hard to integrate into production software\\
\bottomrule
\end{tabularx}
\end{center}

\section{Further reading}
\begin{itemize}
\item S.~Wolfram, \emph{An Elementary Introduction to the Wolfram Language}, 3rd ed.,
  Wolfram Media, 2023 (free online at \url{https://www.wolfram.com/language/elementary-introduction/}).
\item The Wolfram Language documentation centre, \url{https://reference.wolfram.com}.
\item The SymPy tutorial, \url{https://docs.sympy.org}.
\end{itemize}

\begin{exercises}
\item Use Wolfram|Alpha or Mathematica to find the derivative and the indefinite integral
  of $x^2\ln x$. Check the derivative by hand.
\item Solve the drag ODE symbolically with \code{DSolve} (or with SymPy's \code{dsolve})
  and confirm equation~\eqref{eq:dragexact}.
\item The deflection curve of a simply supported beam with a uniform load $q$ satisfies
  $EI\,w''''(x) = q$ with $w(0) = w(L) = 0$ and $w''(0) = w''(L) = 0$. Solve this with
  \code{DSolve} and find the maximum deflection. Compare with $5qL^4/(384EI)$.
\item Use \code{Series} to find the first three terms of the Taylor series of
  $\sqrt{1+x}$ about $x = 0$. Where in engineering is this approximation used?
\item Write the Newton iteration for $x^3 - 2x - 5 = 0$ using \code{NestList} and a pure
  function. How many steps are needed from $x_0 = 2$?
\item Repeat exercises 1 and 2 with SymPy in \py. Which do you find easier?
\end{exercises}
